\subsection{The chemical coordinates of the walls}\label{sec:walls}

Where, chemically, does the curvature mass sit? The eleven
boundary-crossing sweeps of \S\ref{sec:m1} answer directly: each
sweep's per-reaction curvature masses, normalized by that sweep's
own total, aggregate to a per-reaction distribution over the
network, and the distribution is strikingly concentrated. Fifty
reactions --- $1.8\%$ of the network --- carry $85.9\%$ of all
curvature mass, and twenty-three of the twenty-five heaviest
reactions participate in every one of the eleven sweeps. The
heaviest are the enzymes of central carbon metabolism:
glucose-6-phosphate isomerase ($3.7\%$ of all mass), ATP synthase
($3.5\%$), phosphofructokinase and fructose-bisphosphate aldolase
($3.3\%$ each), and the pentose-phosphate entry block
(glucose-6-phosphate dehydrogenase,
6-phosphogluconolactonase, phosphogluconate dehydrogenase; $3.2\%$
each).

Rolling the per-reaction mass up to the metabolites those
reactions share gives the wall coordinates of
Table~\ref{tab:walls}: the forks of central carbon. The
hexose-phosphate pool, the triose-phosphate pool, the PEP/pyruvate
energy fork, the anaplerotic TCA forks, and coenzyme A are the
heaviest coordinates; together the top twenty branch metabolites
carry $59.7\%$ of the total mass. Including the energy and redox
currency layer (ATP, NAD(P)(H), ubiquinone, protons) brings
branch- and currency-adjacent reactions to $98.6\%$ of the mass.
The geometry of rerouting lives, almost entirely, at the chemical
coupling layer --- exactly where one metabolite feeds several
enzymes.

Two honesty notes bound the claim. First, the reaction-level
enrichment is a concentration effect, not a distributional shift:
$46.3\%$ of reactions are branch-adjacent and carry $60.9\%$ of
the mass (permutation $p = 0.017$; under a tighter branch-degree
threshold, $28.1\%$ of reactions carry $50.6\%$, $p = 0.002$),
while the rank AUC is $\sim 0.5$ --- branch degree alone does not
predict a reaction's mass. The strong statement is the
concentration one: the mass that exists attaches to branch
chemistry, and a handful of forks carry most of it. Second, of the
fifty heaviest reactions, $56\%$ are branch-adjacent against the
$46.3\%$ baseline: the concentration is on central branch-point
chemistry, not on branching as such.

\begin{table}[t]
\centering
\caption{Wall coordinates: the ten branch metabolites carrying the
most curvature mass over the eleven boundary-crossing M1 sweeps
(each sweep normalized by its own total mass; each reaction's mass
split over the branch metabolites it connects). The full top-twenty
list carries $59.7\%$ of all curvature mass. Degree is the number
of reactions the metabolite participates in.}
\label{tab:walls}
\small
\begin{tabular}{llrl}
\toprule
metabolite & role & degree & mass share \\
\midrule
fructose-6-phosphate & carbon branch point & 18 & $8.1\%$ \\
glyceraldehyde-3-phosphate & carbon branch point & 15 & $6.9\%$ \\
glucose-6-phosphate & carbon branch point & 14 & $5.5\%$ \\
dihydroxyacetone phosphate & carbon branch point & 16 & $4.9\%$ \\
phosphoenolpyruvate & carbon branch point & 34 & $4.5\%$ \\
2-oxoglutarate & nitrogen branch point & 30 & $3.6\%$ \\
pyruvate & carbon branch point & 72 & $3.6\%$ \\
coenzyme A & acyl-transfer hub & 80 & $3.6\%$ \\
oxaloacetate & carbon branch point & 13 & $2.1\%$ \\
succinate & carbon branch point & 28 & $2.1\%$ \\
\bottomrule
\end{tabular}
\end{table}

The table converts the discrete measure into a wet-lab observable:
targeted metabolomics of roughly twenty pools during a nutrient
transition is a direct experimental readout of the theory's
central object. The prediction is developed in
\S\ref{sec:memory}, where the same pools reappear as the candidate
substrate of metabolic memory.

\subsection{A conserved interior architecture}\label{sec:interior}

A natural expectation is that the walls sit at the network's
boundary with its environment --- at the exchange and transport
reactions, where nutrients enter and leave. The sweeps say
otherwise. Exchange reactions are $12.4\%$ of the network but
carry $3.9\%$ of the curvature mass, and $92.6\%$ of them stay
flat across every sweep; transport reactions are $32.7\%$ and
carry $20.7\%$ ($90.4\%$ flat); internal reactions are $54.9\%$ of
the network and carry $75.4\%$ of the mass ($74.3\%$ flat;
permutation $p = 0.0012$). The interface with the environment is
flat; the walls sit on internal chemistry. The split is stable
across sweep families --- nutrient sweeps place $74.1\%$ of the
mass on internal reactions, enzyme-knockdown sweeps $75.7\%$ --- so
it is not an artifact of perturbing uptake bounds rather than
internal capacities.

The subsystem rollup names the walls: glycolysis and
gluconeogenesis ($23.4\%$ of all curvature mass), the pentose
phosphate pathway ($19.7\%$), the citric acid cycle ($13.9\%$),
oxidative phosphorylation ($11.6\%$), and anaplerosis ($4.7\%$)
--- $73.3\%$ of the mass in five central subsystems, with
branch-point chemistry carrying $74.4\%$ of the internal mass.
The iJO1366 glucose sweep completes the picture from the other
side: its total second-difference mass is $6.5 \times 10^{-9}$
across all $2{,}583$ reactions --- an entire trajectory traversing
a single chamber, every kink an anchor corner of the trajectory
design --- while its active-reaction support moves between $421$
and $440$ reactions, bracketing the endpoint census of $438$
(\S\ref{sec:e22}).

The architectural reading inverts the interface picture. The walls
of the flux geometry sit at the conserved, internal, branch-point
chemistry of central carbon metabolism; the exchange interface is
flat. The species-specific layer is then not the wall chemistry
--- the forks of central carbon are shared across biology --- but
its regulon curation (\S\ref{sec:anatomy}): which regulators price
crossings of the same walls. This offers a mechanism for a
familiar observation, the deep conservation of central branch
chemistry alongside the divergence of its transcriptional wiring:
the chemistry is where the geometry has to be; the regulation is
how each species prices crossings of it. A cross-species
prediction follows, partially checkable with the machinery
already deposited: the branch-metabolite concentration of
curvature mass should be conserved across reconstructions, while
the regulon membership of the same wall-adjacent genes diverges.
